% =========================================================
% Proof Stub: Finite Alteration Does Not Affect Limit
% Source: volume-iii/analysis/sequences/notes/sequences/notes-k-tails.tex
% Mode: standard
% =========================================================

\clearpage
\phantomsection
\hypertarget{proof-RA-ABB-C-finite-alteration}{}

\subsubsection[Finite Alteration Does Not Affect Limit]{Proof --- Finite Alteration Does Not Affect Limit}

\bigskip

\noindent
\textbf{Source.}
See \texttt{notes-k-tails.tex}.

\vspace{0.75em}

\noindent
\textbf{Goal.}
Removing or altering finitely many terms of a convergent sequence does not change its limit.

\vspace{0.75em}

\noindent
\textbf{Logical form.}
\[
  (x_n) \to L,\; (y_n) \text{ differs from } (x_n) \text{ in finitely many terms} \Rightarrow y_n \to L.
\]

\vspace{0.75em}

\noindent
\textbf{Proof strategy.}
This is the Finite Modification Theorem stated as a corollary about convergence.

\vspace{0.75em}

\noindent
\textbf{Proof.}
\begin{proof}
\Claim Removing or altering finitely many terms of a convergent sequence does not change its limit.

\Given 

\Goal 

\Strategy This is the Finite Modification Theorem stated as a corollary about convergence.

\medskip

% --- cite Finite Modification Theorem directly ---

\end{proof}

\begin{remark*}[Proof shape]
Immediate consequence of the Finite Modification Theorem.
\end{remark*}

\begin{remark*}[Dependencies]
\begin{itemize}
  \item Finite Modification Theorem.
\end{itemize}
\end{remark*}
